\begin{afterword}
\summaryhead

The post answers the commenter poke, who said that forming good hypotheses is a skill scientists learn in apprenticeship. The author did not go through that apprenticeship, but sees eminent scientists making the young author's mistakes. Roger Penrose still holds that consciousness is caused by quantum gravity, an answer that leaves consciousness as mysterious as before; the author doubts that Penrose, or Niels Bohr, was ever warned against such answers. The author restates four danger signs of mysterious answers and says no one taught them. No one, the author says, told the physicists who reject many-worlds about the formal definition of Occam's razor. Great mentors may pass on skills they cannot describe, but that is not standard science.

Scientists who describe the future in detail were not taught the conjunction fallacy, and whole subfields grow around vague words such as ``emergence.'' Traditional rationality teaches people to give up a theory under crushing evidence; professional training adds frequentist statistics and standard techniques, and asking more would stop science. As Nick Tarleton said, the lax social standard, taken as a private one, lets people believe too much. A famous mentor may share rare advice, but there is no secret tradition of precise reasoning: half of all scientists ``still believe they believe'' in God.

\responsehead

The title asks an empirical question: what are scientists taught? The post answers it from the author's impressions, with hedges throughout (``I expect,'' ``I doubt,'' ``as far as I can tell''). It offers cases the author judges to be mistakes and reports seeing no sign of a warning. It gives no survey of training.

The evidence that exists, none of it cited, mostly supports the post's general claim. Tversky and Kahneman found in 1971 that professional psychologists had ``erroneous intuitions about the laws of chance.'' Feynman told Caltech's graduates in 1974 that the discipline of not fooling oneself was in no course Feynman knew of: ``We just hope you've caught on by osmosis.'' Zuckerman's study of Nobel laureates found that masters taught ``less by precept than by example.'' Kuhn argued in 1959 that scientific training is a narrow initiation into a tradition, and that this is part of why normal science works, which is close to the post's claim that asking more would stop science.

Where the post goes beyond this evidence, it takes a side without naming the other. It treats mentors' failure to put their guidance into words as a sign that ``these things'' are poorly understood. Michael Polanyi's view, close to poke's, is that the art of research can be passed on ``only by example from master to apprentice'' because no prescription for it exists. The post answers poke without addressing that view. And its two examples of ``rare personal secrets'' had been published in similar words, by Hamming in 1986 and by Feynman in 1974. That does not fit its suggestion that great mentors have not put their guidance into words.

Several cases depend on disputed or unshown premises. Penrose reached the view by an argument from Gödel's theorem that the post does not mention. The many-worlds case depends on a verdict assessed in the notes on ``Many Worlds, One Best Guess.'' The ``exact rational probability'' with no room for whims is one school's view. The figure for scientists who believe in God is roughly that of American surveys, but ``believe they believe'' is asserted without evidence.

In short: an answer from impressions that the evidence it does not cite (Tversky and Kahneman, Feynman, Zuckerman, Kuhn) largely supports, though it presents as rare secrets advice that Hamming and Feynman had published in similar words, and passes over Polanyi's view that such skill cannot be put into rules.
\end{afterword}
